\phantomsection\label{fig:vol01:western-zhou:muye}
\begin{center}
\begin{tikzpicture}[x=.9cm,y=.9cm,font=\small]
  \fill[paper] (0,0) rectangle (10.2,5.5);
  \draw[blue!55,line width=1pt] (.4,4.5) .. controls (2.5,4.2) and (5.3,4.8) .. (9.8,4.2);
  \node[text=blue!70!black,above] at (5.0,4.45) {黄河};
  \draw[blue!45,line width=.75pt] (6.8,.6) .. controls (7.1,1.9) and (7.2,2.7) .. (7.1,3.8);
  \node[text=blue!70!black,right] at (7.2,2.3) {淇水方向};

  \fill[dynasty!23] (.8,1.0) -- (4.0,.9) -- (4.4,3.5) -- (1.3,3.7) -- cycle;
  \node[align=center] at (2.5,2.3) {周方\\自西而来};
  \fill[green!20] (7.4,1.5) -- (9.7,1.4) -- (9.8,3.6) -- (7.6,3.8) -- cycle;
  \node[align=center] at (8.6,2.65) {商都\\朝歌方向};
  \fill[timeline] (5.5,3.1) circle (.14);
  \node[above,align=center] at (5.5,3.2) {牧野\\战场区域};
  \draw[-{Stealth[length=2mm]},ultra thick,color=dynasty] (3.6,2.7) -- (5.35,3.05);
  \draw[-{Stealth[length=2mm]},thick,color=dynasty] (5.65,3.05) -- (7.25,2.8);
  \node[text=dynasty,below] at (4.6,2.65) {周方东进};
\end{tikzpicture}

\smallskip
\small\textit{图  牧野战事空间示意：突出周方由西向东进入商核心区、牧野与朝歌方向的关系；战场范围、行军路线与双方部署均为约略表达。}
\end{center}
